\documentclass[11pt]{article}
\usepackage[margin=1in]{geometry}
\usepackage{amsmath,amssymb,amsthm}
\usepackage{graphicx}
\usepackage[hidelinks]{hyperref}

\newtheorem{theorem}{Theorem}

\newcommand{\D}{\mathcal D}
\newcommand{\Hh}{\mathcal H}
\newcommand{\tr}{\operatorname{tr}}
\newcommand{\supp}{\operatorname{supp}}

\title{A \texorpdfstring{$C^1$}{C1} Tangent Cone for Density Matrices}
\author{FA/Banach run \texttt{fa\_banach\_001}}
\date{2026-07-18}

\begin{document}
\maketitle

\section*{Status}

This packet gives a likely valid full solution to the regularity question
raised in Cai--Govindarajan--Junge, arXiv:2504.04883.  The source defines the
tangent cone to the finite-dimensional density-matrix state space using
one-sided $C^2$ curves, then says that a possible future direction is weakening
the analytic condition to $C^1$.  The result below proves that this weakening
is possible and that the cone is unchanged.

\begin{figure}[h]
\centering
\includegraphics[width=\textwidth]{figures/open_problem_crop.png}
\caption{Source page 14: future-work sentence and the existing $C^2$ tangent cone theorem.}
\end{figure}

\section*{Theorem}

Let $\Hh$ be finite-dimensional and let
\[
  \D(\Hh)=\{\rho\in B(\Hh): \rho\ge 0,\ \tr\rho=1\}.
\]
For $\rho\in \D(\Hh)$, define the one-sided $C^1$ tangent cone
\[
T^{(1)}_\rho \D(\Hh)
  = \left\{\gamma'(0):
      \gamma\in C^1([0,\varepsilon),\D(\Hh)),\ \gamma(0)=\rho
    \right\}.
\]
If $s$ is the orthogonal projection onto $\supp\rho$, and $q=1-s$, then
\[
T^{(1)}_\rho \D(\Hh)
 =
\{x\in B(\Hh): x=x^*,\ \tr x=0,\ qxq\ge 0\}.
\]
In particular, the $C^1$ tangent cone agrees with the $C^2$ tangent cone of
arXiv:2504.04883, Theorem 3.3.  Consequently the source's identity
\[
T_\rho^+\D(\Hh)=\{L(\rho): L\text{ is Lindbladian}\}
\]
remains true when $T_\rho^+\D(\Hh)$ is defined using $C^1$ curves.

\section*{Idea of the Proof}

A $C^1$ path has the expansion $\gamma(t)=\rho+t x+o(t)$.  On the kernel of
$\rho$, positivity gives $q\gamma(t)q=t\,qxq+o(t)\ge 0$.  Dividing by $t$ and
taking the limit gives $qxq\ge 0$.  Conversely, every $x$ satisfying this
condition can be realized by a quadratic correction.  The correction absorbs
the off-diagonal block that prevents the linear path $\rho+t x$ from staying
positive at boundary points.

\section*{Proof}

First suppose $x\in T^{(1)}_\rho \D(\Hh)$.  Then there is a one-sided
$C^1$ curve $\gamma:[0,\varepsilon)\to \D(\Hh)$ with $\gamma(0)=\rho$ and
$\gamma'(0)=x$.  Since every $\gamma(t)$ is self-adjoint and has trace one,
$x=x^*$ and $\tr x=0$.  Also
\[
  \gamma(t)=\rho+t x+r(t),\qquad \|r(t)\|=o(t).
\]
Compress to $q\Hh=\ker\rho$.  Since $q\rho q=0$ and $\gamma(t)\ge 0$,
\[
  0\le q\gamma(t)q=t\,qxq+q r(t)q .
\]
After division by $t>0$ and passage to the limit, closedness of the positive
semidefinite cone gives $qxq\ge 0$.

Conversely, assume $x=x^*$, $\tr x=0$, and $qxq\ge 0$.  If $q=0$, then $\rho$
is positive definite and $\rho+t x\in\D(\Hh)$ for all sufficiently small
$t\ge 0$, so $x$ is realized by a $C^\infty$ path.

It remains to treat the boundary case.  Write matrices relative to
$s\Hh\oplus q\Hh$:
\[
  \rho=
  \begin{pmatrix} R&0\\0&0\end{pmatrix},
  \qquad
  x=
  \begin{pmatrix} A&B\\B^*&C\end{pmatrix},
\]
where $R>0$ on $s\Hh$, $C\ge 0$ on $q\Hh$, and
$\tr A+\tr C=0$.  Choose $D\ge 0$ on $q\Hh$ so large that
\[
  D\ge 2B^*R^{-1}B .
\]
For small $t\ge 0$, define
\[
  \rho_t=
  \begin{pmatrix}
  R+tA-t^2(\tr D)R & tB\\
  tB^* & tC+t^2D
  \end{pmatrix}.
\]
This curve is polynomial in $t$, satisfies $\rho_0=\rho$, and has derivative
$x$ at $0$.  Its trace is
\[
  1+t(\tr A+\tr C)-t^2(\tr D)\tr R+t^2\tr D=1,
\]
because $\tr R=\tr\rho=1$.

For all sufficiently small $t$, the upper-left block
\[
R_t=R+tA-t^2(\tr D)R
\]
is positive definite and satisfies $R_t^{-1}\le 2R^{-1}$ in the Loewner order.
The Schur complement of $R_t$ in $\rho_t$ is therefore
\[
  tC+t^2D-t^2B^*R_t^{-1}B
  \ge
  tC+t^2(D-2B^*R^{-1}B)
  \ge 0.
\]
Thus $\rho_t\ge 0$ for all sufficiently small $t\ge 0$.  We have constructed a
$C^2$ path in $\D(\Hh)$ with derivative $x$, hence certainly a $C^1$ path.

Combining the two inclusions proves the formula for $T^{(1)}_\rho\D(\Hh)$.
The final Lindbladian assertion follows by comparing this formula with
Theorem 3.8 of the source, which identifies the same cone with
$\{L(\rho):L\text{ is Lindbladian}\}$.

\section*{Verification Notes}

The only point where boundary singularity enters is the Schur-complement
estimate.  The $t^2D$ term dominates the order-$t^2$ off-diagonal obstruction,
while the $-t^2(\tr D)R$ term preserves trace without destroying positivity of
the support block.  No compactness or infinite-dimensional closure issue is
involved.

\section*{Novelty Check and Limitations}

On 2026-07-18, the run indexes and web searches were checked for arXiv:2504.04883,
``How Far do Lindbladians Go'', ``C1 tangent cone'', ``density matrices tangent
cone support projection'', ``positive semidefinite C1 tangent cone'', and
``Lindbladian tangent cone''.  The search found the source paper itself and
general convex-analysis material, but no exact later answer to the source's
$C^1$ weakening direction.

This packet does not solve the broader lifting/retract problem for selecting
Lindbladian generators continuously, smoothly, or analytically along arbitrary
paths.  It only proves that the tangent-cone definition can be weakened from
$C^2$ to $C^1$.

\begin{thebibliography}{9}
\bibitem{CGJ}
J. Cai, A. Govindarajan, and M. Junge,
\emph{How Far do Lindbladians Go?},
arXiv:2504.04883.
\end{thebibliography}

\end{document}
